\section{Related Work}
\label{sec:related_work}

The evolution of attention as a concept has traversed multiple domains, from its inception in neural machine translation to its central role in the digital economy. In this section, we review the literature that bridges these fields and positions our work within this broader context.

\para{Attention in Generative AI}
The "Attention is All You Need" paradigm \cite{vaswani2017attention} replaced recurrent neural networks with the self-attention mechanism, enabling massive parallelization and the capture of long-range dependencies in sequence modeling. This approach has since been extended to various modalities, including computer vision \cite{rombach2022high} and large-scale language understanding \cite{devlin2018bert}. While these models excel at processing complex, high-dimensional data, the underlying selective weighting over a limited information budget remains their core innovation. Our work builds on this technical foundation but extends its application to industrial Internet tasks like sequential recommendation and CTR prediction.

\para{Attention in Recommendation and CTR Prediction}
The integration of attention mechanisms into recommendation systems and advertising has led to significant advancements. Models like SASRec \cite{kang2018self} and BERT4Rec \cite{sun2019bert4rec} have demonstrated the power of self-attention for modeling sequential user behaviors. Similarly, in the domain of Click-Through Rate (CTR) prediction, the Deep Interest Network (DIN) \cite{zhou2018deep} uses attention to dynamically weight historical user behaviors relative to a target ad. Recent developments such as RecGRELA \cite{hu2025gated} and PANTHER \cite{panther2025panther} continue to push the boundaries by using linear and rotary-enhanced attention for industrial-scale sequential behavior modeling. Unlike these works that focus on specific silos, our research provides a unified empirical study comparing attention and non-attention baselines across multiple tasks under a consistent statistical framework.

\para{The Attention Economy}
The concept of the "Attention Economy" recognizes that in a world with an abundance of information, human attention becomes the scarcest resource. Theoretical models have proposed mathematical frameworks for the co-evolution of content quality and audience selectivity under limited attention capacity \cite{chujyo2026attention}. Further research has introduced frameworks for measuring "attentional agency" on digital platforms, distinguishing between "push" and "pull" value \cite{wojtowicz2024push}. Our work complements these theoretical perspectives by providing empirical evidence from AI model behavior, linking internal mathematical properties (attention entropy) to the effectiveness of information filtering in digital systems. This connection bridges the gap between the technical mechanism of information mediation and the strategic objectives of the broader digital ecosystem \cite{castaldo2022online}.
